\section{Limited-Memory BFGS and Stochastic Limitations}
\label{sec:lbfgs}

While the BFGS algorithm successfully achieves second-order convergence using only first-order gradients, it fundamentally requires maintaining and updating the dense $d \times d$ inverse Hessian approximation $B_t$. For modern deep learning models where $d$ exceeds tens of millions of parameters, allocating $O(d^2)$ memory is strictly impossible.

Limited-Memory BFGS (L-BFGS) \cite{liu1989limited} resolves this by entirely abandoning the explicit matrix $B_t$. Instead of storing the full matrix, L-BFGS reconstructs the matrix-vector product $B_t \nabla f(w^{(t)})$ dynamically at every step using only a small, truncated history of the most recent parameter and gradient updates.

\subsection{The Two-Loop Recursion}
By unrolling the exact BFGS update equation, one can prove that the current inverse Hessian approximation $B_t$ is entirely defined by an initial guess $B_t^{(0)}$ and the sequence of historical parameter steps $s_k$ and gradient differences $y_k$. 

L-BFGS introduces a memory budget $m$ (typically $5 \le m \le 20$). It stores only the $m$ most recent pairs $\{(s_i, y_i)\}_{i=t-m}^{t-1}$. To compute the optimal search direction $p_t = -B_t g_t$, the algorithm applies a mathematically exact procedure known as the \textbf{Two-Loop Recursion}.

\begin{algorithm}[H]
\caption{L-BFGS Two-Loop Recursion}
\begin{algorithmic}[1]
\Require Gradient $g_t = \nabla f(w^{(t)})$
\Require History $\{(s_i, y_i)\}_{i=t-m}^{t-1}$
\State $q \gets g_t$
\For{$i = t-1$ down to $t-m$} \Comment{Backward Pass}
    \State $\rho_i \gets \frac{1}{y_i^T s_i}$
    \State $\alpha_i \gets \rho_i s_i^T q$
    \State $q \gets q - \alpha_i y_i$
\EndFor
\State $\gamma_t \gets \frac{s_{t-1}^T y_{t-1}}{y_{t-1}^T y_{t-1}}$ \Comment{Initial Hessian Scaling}
\State $r \gets \gamma_t q$
\For{$i = t-m$ to $t-1$} \Comment{Forward Pass}
    \State $\beta \gets \rho_i y_i^T r$
    \State $r \gets r + s_i (\alpha_i - \beta)$
\EndFor
\State \Return $p_t = -r$
\end{algorithmic}
\end{algorithm}

The scalar factor $\gamma_t$ serves as the initial Hessian approximation $B_t^{(0)} = \gamma_t I$, effectively estimating the underlying scale of the curvature before applying the $m$ rank-two updates. 

\textbf{Computational Complexity:} Because the algorithm relies exclusively on sequential vector inner products, computing the optimal search direction requires exactly $O(md)$ operations and $O(md)$ memory. Since $m$ is a small constant, the time and space complexity scale linearly with respect to the parameter count, rendering the update step as computationally efficient as standard Gradient Descent.

\begin{figure}[htbp]
    \centering
    \includegraphics[width=0.85\textwidth]{floats/lbfgs_memory.pdf}
    \caption{L-BFGS convergence on a dense, highly correlated 100-dimensional quadratic objective ($\kappa = 1000$). Standard GD decays infinitely slowly. A memory budget of $m=1$ captures only scalar scaling, struggling against dense cross-correlations. A budget of $m=20$ maps out the high-dimensional correlated curvature flawlessly, achieving exponential second-order convergence without ever forming a $100 \times 100$ matrix.}
    \label{fig:lbfgs_memory}
\end{figure}

\subsection{The Stochastic Limitation of Quasi-Newton Methods}
Given its linear time complexity and massive second-order acceleration, L-BFGS is universally regarded as the optimal algorithm for deterministic, full-batch optimization. However, it is fundamentally incompatible with the training of modern deep neural networks via mini-batching. The reason lies strictly in the mathematical geometry of the Secant Condition.

Recall the gradient difference vector:
\begin{equation}
    y_t = \nabla f(w^{(t+1)}) - \nabla f(w^{(t)})
\end{equation}
For $y_t$ to accurately measure the physical curvature of the loss surface, the two gradients must be evaluated on the \textit{exact same objective function}. 

In Stochastic Gradient Descent (mini-batching), the objective function changes at every step because the sampled data batch changes. If $g_{t+1}$ is evaluated on Batch $A$ and $g_t$ is evaluated on Batch $B$, the difference $y_t$ is hopelessly corrupted by sampling variance. The secant condition is violated, the curvature estimate becomes physically meaningless, and the implicit inverse Hessian $B_t$ loses its positive definiteness, causing the algorithm's line search to fail or the optimizer to diverge entirely.

Consequently, while Adam and RMSProp successfully dampen stochastic noise using diagonal variance accumulators, quasi-Newton methods remain strictly constrained to domains where gradients are fully deterministic or where noise is rigidly controlled via massive batch sizes.
